\newcommand{\docshorttitle}{OMNIBUS SUPPL.\ P\&A --- BLACK-LETTER CONTRADICTIONS (CGC-25-631801)}
\newcommand{\doctitleblock}{PLAINTIFF'S OMNIBUS SUPPL.\ P\&A --- BLACK-LETTER CONTRADICTIONS (MULTI-MOTION DOCKET)}
\newcommand{\docsubblock}{(Cal.\ Rules of Court, rules 3.1113, 3.1306(c); supplemental only; cross-noticed hearings on caption.)}
% Pleading-paper preamble for APRIL-24-2026 cover sheets and oral-argument
% outlines (CGC-25-631801 and CGC-25-631802). Matches house style from
% APRIL-23-SAC-DEFENSE/court-pdfs/REPLY-SAC-PREAMBLE.tex, simplified for
% single-to-three-page artifacts that do not require TOC/TOA infrastructure.
%
% Cross-hearing caption (CAPTION-801/802.tex): use \CaptionDocTitleBlock,
% \CaptionDocSubBlock, and the crosshearingsschedule environment (defined below)
% so long \doctitleblock text and hearing dates stay inside the 3in minipage.
\documentclass{article}
\usepackage[utf8]{inputenc}
\usepackage{graphicx}
\usepackage{geometry}
    \geometry{top=1in, bottom=2.2in, left=1.0in, right=1.0in, textheight=336pt}
    \setlength{\footskip}{70pt}
\usepackage{ulem}
\usepackage{fancyhdr}
    \renewcommand{\headrulewidth}{0pt}
    \renewcommand{\footrulewidth}{0.4pt}
\usepackage{parskip}
    \setlength{\parskip}{0mm}
\usepackage{quoting}
    \quotingsetup{leftmargin=0.5in, rightmargin=0.5in, vskip=-1.5mm}
    \makeatletter
        \g@addto@macro\quoting\singlespacing
        \g@addto@macro\quoting{\vspace{-2mm}}
    \makeatother
\usepackage{mathptmx}
\usepackage{amssymb}
\renewcommand{\rmdefault}{ptm}
\renewcommand{\normalsize}{\fontsize{12}{12}\selectfont}
\newcommand*{\ptm}{\fontfamily{ptm} \selectfont \fontsize{12}{12}\selectfont}
\usepackage{setspace}
    \doublespacing
\newcommand{\tab}{\hspace*{.0in}}
\newenvironment{tightcenter}{%
    \setlength\topsep{0pt}
    \setlength\parskip{0pt}
    \begin{center}
}{%
    \end{center}
}
\usepackage{tikz}
\usetikzlibrary{calc}
% Pleading-paper vertical double rule at the left margin.
\AddToHook{shipout/background}{%
    \begin{tikzpicture}[remember picture, overlay]
        \draw[line width=0.5pt] ($(current page.north west) + (0.75in,0)$) -- ($(current page.south west) + (0.75in,0)$);
        \draw[line width=0.5pt] ($(current page.north west) + (0.80in,0)$) -- ($(current page.south west) + (0.80in,0)$);
    \end{tikzpicture}%
}
\usepackage[left, pagewise]{lineno}
\setlength{\linenumbersep}{0.35in}
\renewcommand{\linenumberfont}{\normalfont\small}
\usepackage{titlesec}
\usepackage[hidelinks]{hyperref}
\usepackage{bookmark}
\usepackage{textcomp}
\usepackage{longtable}
\usepackage{booktabs}
\usepackage{array}
\usepackage{multirow}
\usepackage{calc}
% --- Cross-hearing caption cover (3in minipage): compact title + 3-column hearing table ---
% Use on first page only; keeps long \doctitleblock / \docsubblock on one sheet with party caption.
\newcommand{\CaptionDocTitleBlock}{%
  \begingroup
  \fontsize{10}{11.5}\selectfont
  \bfseries
  \raggedright
  \setlength{\parskip}{0pt}%
  \noindent\doctitleblock\par
  \endgroup
}
\newcommand{\CaptionDocSubBlock}{%
  \begingroup
  \footnotesize
  \raggedright
  \setlength{\parskip}{0pt}%
  \noindent\docsubblock\par
  \endgroup
}
% #1 = department number only (e.g.\ 301 or 302) for the line ``Dept.\ #1''.
\newenvironment{crosshearingsschedule}[1]{%
  \par\begingroup
  \footnotesize
  \setlength{\tabcolsep}{2.2pt}%
  \renewcommand{\arraystretch}{1.04}%
  \setlength{\parskip}{0pt}%
  \setlength{\parindent}{0pt}%
  \noindent\textbf{Cross-noticed for pending defense hearings (this case, Dept.\ #1):}\par\vspace{2pt}%
  \begin{tabular}{@{}>{\raggedright\arraybackslash}p{0.11\linewidth}>{\raggedright\arraybackslash}p{0.34\linewidth}>{\raggedright\arraybackslash}p{0.53\linewidth}@{}}%
  \textbf{\#} & \textbf{Motion} & \textbf{Date / dept / judge} \\ \hline
}{%
  \end{tabular}\par\endgroup\smallskip
}
% Pandoc pipe-table column widths use \real{#}.
\newcommand{\real}[1]{\pgfmathparse{#1}\pgfmathresult}
\titleformat{\section}{\fontsize{13}{13}\selectfont\normalfont\bfseries\uppercase}{}{0em}{}
\titleformat{\subsection}{\fontsize{12}{12}\selectfont\normalfont\bfseries}{}{0em}{}
\titleformat{\subsubsection}{\normalfont\bfseries}{}{0em}{}
\titlespacing*{\section}{0pt}{6pt}{2pt}
\titlespacing*{\subsection}{0pt}{4pt}{2pt}
\titlespacing*{\subsubsection}{0pt}{2pt}{2pt}
\newcommand{\calrules}[1]{Cal. Rules of Court, rule #1}
\newcommand{\ccp}[1]{Code Civ. Proc., \textsection~#1}
\newcommand{\evidcode}[1]{Evid. Code, \textsection~#1}
\providecommand{\tightlist}{%
  \setlength{\itemsep}{0pt}\setlength{\parskip}{0pt}}
% Footer: page N of total + document short-title.
\usepackage{lastpage}
\pagestyle{fancy}
\fancyhf{}
\fancyfoot[C]{\thepage{} of \pageref{LastPage} \\[1pt] \footnotesize \docshorttitle}
\fancyfoot[R]{}

\begin{document}
% Cross-hearing caption for CGC-25-631801 (Dept. 301, Hon. Van Aken).
\nolinenumbers
\begin{singlespace}
\noindent Franciscus Dylan Rosario \\
7624 Melody Drive \\
Rohnert Park, CA 94928 \\
Tel: 650.488.7510 \\
E: soltrinox@gmail.com \\
Plaintiff, In Pro Per
\end{singlespace}
\vspace*{9mm}
\begin{tightcenter}
\textbf{SUPERIOR COURT OF THE STATE OF CALIFORNIA} \\
\textbf{COUNTY OF SAN FRANCISCO}
\end{tightcenter}
\vspace*{3mm}
\begin{minipage}[t]{3in}
    \raggedright
    \textbf{FRANCISCUS DYLAN ROSARIO,} \\ \textbf{PLAINTIFF}, \\
    \hspace{1cm} v. \\
    \small
    Subhi Abdelhalim; CSAA Insurance Exchange; Phillips, Spallas \& Angstadt LLP; Priya D. Navaratnasingham; Alberto Reyna; Michael R. Chambers; Carbone, Smith \& Koyama; and Does 1 through 50, inclusive;\\
    \normalsize
    \textbf{DEFENDANTS}
    \makebox[3in]{\hrulefill}
\end{minipage}%
\hspace{3mm}%
\begin{minipage}[t]{0.5pt}
    \vspace{0pt}
    \rule{0.5pt}{2.42in}
\end{minipage}%
\hspace{3mm}%
\begin{minipage}[t]{3in}
    \raggedright
    \noindent\textbf{Case No.: CGC-25-631801}\par\smallskip
    \CaptionDocTitleBlock\smallskip
    \CaptionDocSubBlock\smallskip
    \begin{singlespace}
    \setlength{\parindent}{0pt}
    \begin{crosshearingsschedule}{301}
    (1) & MFL / leave (SAC) & \textbf{Apr.\ 30, 2026} \\
    (2) & Anti-SLAPP (\ccp{425.16}) & \textbf{May 12, 2026, 9:30 a.m.};\ Dept.\ 301;\ Hon.\ Christine Van Aken \\
    (3) & Strike (\ccp{436}) & \textbf{May 22, 2026};\ Dept.\ 301;\ Hon.\ Christine Van Aken \\
    (4) & Demurrer & Hrng.\ date \textbf{[TBD]} (reserved) \\
    \end{crosshearingsschedule}
    \noindent Complaint filed: Dec.\ 23, 2025 \\
    \noindent FAC filed: Feb.\ 25, 2026 \\
    \end{singlespace}
\end{minipage}
\vspace*{8mm}
\newpage
\linenumbers


\section*{I. INTRODUCTION AND VEHICLE}

This omnibus supplemental memorandum is submitted under \calrules{3.1113} and \calrules{3.1306(c)} as a supplement to Plaintiff's previously filed oppositions to Defendants' Anti-SLAPP motion (Apr.\ 9, 2026), demurrer (Apr.\ 13, 2026), \ccp{436} motion (Apr.\ 22, 2026), and opposition to Plaintiff's motion for leave to file a second amended complaint (Apr.\ 23, 2026). It is \textbf{not} a substitute for those oppositions. It maps recurring defense positions to contradicting California authority using a single template: defense position (verbatim cite to the defense extraction record) $\rightarrow$ contradicting authority (pin cite) $\rightarrow$ defect class $\rightarrow$ disposition. Panel-counsel and insurer postures are addressed together where the record supports a unified regulatory and agency framing.

\section*{II. ISSUES PRESENTED (REBUTTED DEFENSE THEORIES)}

\begin{itemize}
\tightlist
\item Whether ``absolute'' litigation-privilege and \ccp{425.16} tropes erase threshold \ccp{425.17(c)} sequencing and Insurance Code / Title 10 duties (\textbf{No}).
\item Whether a third-party claimant is excluded from the customer / regulatory audience protected by \ccp{425.17(c)(2)} and Title 10 (\textbf{No}; see concurrently filed Customer-Doctrine Supplemental).
\item Whether demurrer four-corners analysis may import merits-stage disparagement and collateral-attack rhetoric (\textbf{No}; \textit{City of Cotati v. Cashman} (2002) 29 Cal.4th 69).
\item Whether a \ccp{436} motion may operate as a backdoor anti-SLAPP or demurrer (\textbf{No}; \textit{PH II, Inc.\ v.\ Superior Court} (1995) 33 Cal.App.4th 1680).
\item Whether MFL/SAC opposition may collapse \textit{Park v.\ Board of Trustees} step-one gravamen into step-two probability language (\textbf{No}; \textit{Park} (2017) 2 Cal.5th 1057; \textit{Baral v.\ Schnitt} (2016) 1 Cal.5th 376).
\item Whether panel counsel may disclaim the commercial channel for \ccp{425.17(c)} purposes while performing insurer-directed claim and investigation communications (\textbf{No}; \textit{JAMS, Inc.\ v.\ Superior Court} (2016) 1 Cal.App.5th 984; Civ.\ Code, \S~2330).
\end{itemize}

\section*{III. REGULATORY AND STATUTORY FLOOR}

Insurance Code \S~790.02 frames the Unfair Insurance Practices Act as regulating trade practices in the business of insurance without limiting protection to named insureds alone. Insurance Code \S~790.03(h) enumerates unfair claims settlement practices; twelve subdivisions speak to ``claimants'' or parallel ``insureds or claimants'' phrasing, refuting any first-party-only gloss. Insurance Code \S~790.10 delegates authority for the Fair Claims Settlement Practices Regulations. Ten CCR \S\S~2695.1(a), (d), 2695.2(c), (s), (t), 2695.4(a), 2695.5, and 2695.7(a)--(d) define ``claimant'' to include third parties, impose duties on ``all claims,'' and require diligent investigation and timely written disposition. \ccp{425.17(c)} supplies a categorical commercial-speech sequencing rule (\textit{Simpson Strong-Tie Co.\ v.\ Gore} (2010) 49 Cal.4th 12; \textit{Club Members for an Honest Election v.\ Sierra Club} (2008) 45 Cal.4th 309). Civil Code \S~47(b)(2)--(3) and Bus.\ \& Prof.\ Code, \S~6128(a) mark limits on ``petitioning'' re-characterizations. Insurance Code \S~700 et seq.\ conditions the privilege of licensure on compliance with those duties.

\section*{IV. ANTI-SLAPP MPA (APR.\ 9, 2026) --- ROWS A-1 THROUGH A-9}

Each row: \textbf{Defense position} $\rightarrow$ \textbf{Contradicting authority} $\rightarrow$ \textbf{Defect} $\rightarrow$ \textbf{Disposition}.

\begin{description}
\tightlist
\item[A-1.] ``Early dismissal of unmeritorious claims'' policy rhetoric treated as dispositive of step one. \textit{Club Members}, 45 Cal.4th at pp.\ 316--324; \textit{Simpson Strong-Tie}, 49 Cal.4th at pp.\ 28--32. \textit{Mischaracterized.} Treat \ccp{425.17} as structurally prior.
\item[A-2.] Two-step \ccp{425.16} framework quoted to skip \ccp{425.17} gateway. \textit{Club Members} (sequencing); \ccp{425.17(e)} (no interlocutory appeal on exemption orders). \textit{Omitted.} Address carve-out first.
\item[A-3.] ``Collateral attack'' / ``unsupportable'' labels on the FAC at step two. \textit{City of Cotati v.\ Cashman}, 29 Cal.4th at p.\ 76 (demurrer standards govern pleading attacks); \textit{Navellier v.\ Sletten} (2002) 29 Cal.4th 82, 89--92. \textit{Overbroad.} Separate rhetoric from rule.
\item[A-4.] ``Absolute'' litigation privilege as blanket immunity. \textit{Action Apartment Ass'n v.\ City of Santa Monica} (2007) 41 Cal.4th 1232, 1245--1252 (\S~47(b) coextensive with \ccp{425.16(e)} for certain channels); Civ.\ Code, \S~47(b)(2)--(3) (express exceptions). \textit{Overbroad.}
\item[A-5.] Probability / ``meritless'' characterization at step one. \textit{Park}, 2 Cal.5th at pp.\ 1062--1063; \textit{Baral}, 1 Cal.5th at pp.\ 396--399. \textit{Misinterpreted.}
\item[A-6.] \textit{Flatley v.\ Mauro} (2006) 39 Cal.4th 299 and \textit{Lefebvre v.\ Lefebvre} (2011) 199 Cal.App.4th 696 reserved for illegality / uncontroverted-record channels---not a license for generic disparagement. \textit{Mischaracterized} if invoked beyond their holdings.
\item[A-7.] Denial of customer / commercial-channel framing for insurer speech. \textit{Simpson Strong-Tie}; \textit{JAMS}; \textit{Stewart v.\ Rolling Stone LLC} (2010) 181 Cal.App.4th 664; Title 10, \S\S~2695.2(c), (t). \textit{Misinterpreted.}
\item[A-8.] UCL / fraud predicates mislabeled as pure ``petitioning.'' \textit{Rusheen v.\ Cohen} (2006) 37 Cal.4th 1048, 1061--1064 (distinction between protected petitioning and non-communicative or commercial misconduct giving rise to liability). \textit{Omitted} branch analysis.
\item[A-9.] Panel counsel as ``litigation-only'' actors outside \ccp{425.17(c)}. \textit{JAMS}, 1 Cal.App.5th at pp.\ 994--995 (commercial channel includes lawyers delivering the regulated product); Civ.\ Code, \S~2330 (principal--agent alignment for ostensible authority). \textit{Misinterpreted.}
\end{description}

\section*{V. DEMURRER MPA (APR.\ 13, 2026) --- ROWS D-1 THROUGH D-7}

\begin{description}
\tightlist
\item[D-1.] Four-corners attack imports step-two ``frivolous'' rhetoric. \textit{Cashman}, 29 Cal.4th at p.\ 76. \textit{Overbroad.}
\item[D-2.] Grounds framed as pure law questions while smuggling merits disparagement. \textit{Cashman}, 29 Cal.4th at pp.\ 76--78. \textit{Mischaracterized.}
\item[D-3.] Unfair competition / insurance-practice theories treated as categorically implausible on the face of regulated duties. Ins.\ Code, \S\S~790.03(h)(1), (h)(3), (h)(5); 10 CCR \S~2695.7(d). \textit{Omitted} regulatory floor.
\item[D-4.] ``Collateral attack'' on underlying judgment used to obscure discrete new-duty theories. \textit{Cashman}; \textit{Baral} (leave standards). \textit{Misinterpreted.}
\item[D-5.] Privilege and \S~47(b) used to defeat fraud / investigation-duty allegations at pleading stage. \textit{Action Apartment}; Civ.\ Code, \S~47(b)(3). \textit{Overbroad.}
\item[D-6.] Investigation / denial communications treated as non-actionable ``opinion.'' \textit{Egan v.\ Mutual of Omaha Ins.\ Co.} (1979) 24 Cal.3d 809, 818--819; \textit{Wilson v.\ 21st Century Ins.\ Co.} (2007) 42 Cal.4th 713, 720--721. \textit{Mischaracterized.}
\item[D-7.] Agency / non-delegable flavor: insurer and retained counsel relationship to claim investigation. Civ.\ Code, \S~2330; 10 CCR \S~2695.7(d) (fair and objective investigation); \textit{Engalla v.\ Permanente Medical Group} (1997) 15 Cal.4th 951 (managed-care context for institutional accountability principles). \textit{Omitted} alternative framing.
\end{description}

\section*{VI. CCP 436 MOTION (APR.\ 22, 2026) --- ROWS S-1 THROUGH S-7}

\begin{description}
\tightlist
\item[S-1.] \ccp{436} used to restrike claims already analyzed under demurrer / anti-SLAPP without honest narrowing. \textit{PH II}, 33 Cal.App.4th at pp.\ 1682--1683. \textit{Overbroad.}
\item[S-2.] ``Strategic lawsuit'' rhetoric applied to fraud and UCL theories governed by regulatory standards. \textit{Rusheen}; Ins.\ Code, \S~790.03(h). \textit{Mischaracterized.}
\item[S-3.] Prior ``on the merits'' trial outcome conflated with whether carrier pre-denial conduct was litigated. \textit{Park} (gravamen). \textit{Misinterpreted.}
\item[S-4.] Anti-SLAPP policy blocks imported wholesale. \textit{Club Members}; \textit{PH II}. \textit{Omitted} sequencing.
\item[S-5.] Strike motion as surrogate for fact finding on investigation adequacy. 10 CCR \S~2695.7(d); \textit{Wilson}, 42 Cal.4th at pp.\ 720--721. \textit{Overbroad.}
\item[S-6.] Litigation privilege duplicates used to shield pre-trial representations to tribunal. Civ.\ Code, \S~47(b)(2)--(3); Bus.\ \& Prof.\ Code, \S~6128(a). \textit{Misinterpreted.}
\item[S-7.] Equitable fraud-on-the-court boilerplate in lodged exhibits distinguished from defense-authored misstatements of standard---but defense must not misapply \textit{Flatley} to immunize investigation misrepresentations. \textit{Flatley}, 39 Cal.4th at pp.\ 316--322. \textit{Mischaracterized} if over-read.
\end{description}

\section*{VII. MFL/SAC OPPOSITION (APR.\ 23, 2026) --- ROWS M-1 THROUGH M-4}

\begin{description}
\tightlist
\item[M-1.] Step-two / probability language at step one. \textit{Park}; \textit{Baral}. \textit{Misinterpreted.}
\item[M-2.] Service / procedural objections. \textbf{Conceded} as potentially curable; does not defeat threshold sequencing once cured.
\item[M-3.] Lodging / formatting objections. \textbf{Conceded} as procedural; do not substitute for \textit{Park} gravamen analysis.
\item[M-4.] ``Frivolous'' anti-SLAPP fee threats used as merits proxy. \ccp{425.16}(c)(1); \textit{JKC3H8 v.\ Colton} (2013) 221 Cal.App.4th 468; \textit{Nguyen-Lam v.\ Cao} (2009) 171 Cal.App.4th 871. \textit{Overbroad} if invoked to skip step one.
\end{description}

\section*{VIII. CROSS-CUTTING --- CUSTOMER DOCTRINE AND INSURANCE-CODE TABLE}

Plaintiff incorporates by reference the analysis in the concurrently filed Supplemental Memorandum titled \textit{Customer Doctrine under Code of Civil Procedure section 425.17(c)} (Customer-Doctrine Supplemental, lodged concurrently).

\begin{center}
\footnotesize
\begin{tabular}{@{}p{1.6in}p{4.2in}@{}}
\toprule
\textbf{Provision} & \textbf{Black-letter point for omnibus use} \\
\midrule
Ins.\ Code \S~790.02 & Regulates trade practices in the business of insurance; no insured-only limitation stated. \\
Ins.\ Code \S~790.03(h) & Sixteen unfair practices; ``claimants'' and ``insureds or claimants'' phrasing. \\
Ins.\ Code \S~790.10 & Delegation for Title 10 regulations. \\
Ins.\ Code \S~11580(b)(2) & Duty-to-defend context informs tripartite relationships (notice parallel in RJN if reproduced). \\
Ins.\ Code \S~791.13 & Consumer comparison / unfair method context where pleaded. \\
10 CCR \S~2695.2(c), (t) & ``Claimant'' includes third-party claimant. \\
10 CCR \S~2695.7(b), (d) & Written denial bases; thorough, fair, objective investigation. \\
\ccp{425.17(c)(1)--(2)} & Commercial actor; commercial speech / delivery; triply disjunctive audience. \\
\bottomrule
\end{tabular}
\end{center}

\section*{IX. APPLICATION TO CGC-25-631801}

Defendants include CSAA Insurance Exchange and panel counsel. The omnibus matrix applies across calendars: threshold \ccp{425.17(c)} and Title 10 duties are common spine issues; demurrer and \ccp{436} papers must not import merits disparagement as pleading-stage law; MFL/SAC opposition must not collapse anti-SLAPP steps. Plaintiff preserves all prior opposition arguments and adds only this structured contradiction map.

\section*{X. CONCLUSION AND RELIEF REQUESTED}

For each hearing, Plaintiff requests that the Court (a) treat this filing as supplemental authority under \calrules{3.1306(c)}; (b) reject defense over-reads flagged above; (c) deny or sustain objections consistent with the cited authorities; and (d) take judicial notice as requested in the concurrently amended RJN.

\appendix
\section*{APPENDIX A. DEFENSE POSITION INVENTORY (REPRESENTATIVE ROWS)}

\footnotesize
\noindent\textit{Note: Line numbers refer to extraction.md in each defense bundle; confirm quotations against source.pdf before filing. See Merit Isolation Report on file with counsel.}

\begin{longtable}{@{}p{0.35in}p{0.9in}p{0.55in}p{1.35in}p{1.35in}p{0.45in}@{}}
\toprule
\textbf{ID} & \textbf{Pleading} & \textbf{Line} & \textbf{Quote (trunc.)} & \textbf{Contradicting authority} & \textbf{Class} \\
\midrule
\endfirsthead
\toprule
\textbf{ID} & \textbf{Pleading} & \textbf{Line} & \textbf{Quote (trunc.)} & \textbf{Contradicting authority} & \textbf{Class} \\
\midrule
\endhead
801-A & Anti-SLAPP MPA & 484 & ``The absolute litigation privilege \dots'' & Civ.\ Code, \S~47(b)(2)--(3); \textit{Action Apartment} & Overbroad \\
801-B & Anti-SLAPP MPA & 1713 & ``Unsupportable `Collateral Attack'\,'' & \textit{Cashman}, 29 Cal.4th 69 & Mischaracterized \\
801-C & Demurrer MPA & 1834 & ``frivolous new lawsuit'' & \textit{Cashman}; \ccp{128.7}(b) safe harbor vs.\ rhetoric & Rhetoric \\
801-D & CCP 436 MPA & 676 & ``legally and factually unsupportable theory'' & \textit{Park}; \textit{PH II} & Overbroad \\
801-E & MFL dec-toc & 1128 & ``No single ground has merit'' & \textit{Baral} (demurrer v.\ anti-SLAPP distinct roles) & Misinterpreted \\
801-F & MFL opposition & 129 & ``legally unsupportable motion practice'' & \textit{Park} step separation & Rhetoric \\
\bottomrule
\end{longtable}

\section*{APPENDIX B. AUTHORITIES CITED}

\begin{enumerate}
\tightlist
\item Cal.\ Rules of Court, rules 3.1113, 3.1306(c).
\item Code Civ.\ Proc., \S\S~128.7(b), 425.16, 425.17, 436, 904.1(a)(13).
\item Civ.\ Code, \S\S~47(b), 2330.
\item Bus.\ \& Prof.\ Code, \S~6128(a).
\item Ins.\ Code, \S\S~700 et seq., 790.02, 790.03(h), 790.10, 11580(b)(2), 791.13.
\item Cal.\ Code Regs., tit.\ 10, \S\S~2695.1--2695.7.
\item \textit{Simpson Strong-Tie Co.\ v.\ Gore} (2010) 49 Cal.4th 12.
\item \textit{Club Members for an Honest Election v.\ Sierra Club} (2008) 45 Cal.4th 309.
\item \textit{JAMS, Inc.\ v.\ Superior Court} (2016) 1 Cal.App.5th 984.
\item \textit{Stewart v.\ Rolling Stone LLC} (2010) 181 Cal.App.4th 664.
\item \textit{Flatley v.\ Mauro} (2006) 39 Cal.4th 299; \textit{Lefebvre v.\ Lefebvre} (2011) 199 Cal.App.4th 696.
\item \textit{Action Apartment Ass'n v.\ City of Santa Monica} (2007) 41 Cal.4th 1232.
\item \textit{Rusheen v.\ Cohen} (2006) 37 Cal.4th 1048.
\item \textit{City of Cotati v.\ Cashman} (2002) 29 Cal.4th 69.
\item \textit{Navellier v.\ Sletten} (2002) 29 Cal.4th 82.
\item \textit{Park v.\ Board of Trustees} (2017) 2 Cal.5th 1057.
\item \textit{Baral v.\ Schnitt} (2016) 1 Cal.5th 376.
\item \textit{PH II, Inc.\ v.\ Superior Court} (1995) 33 Cal.App.4th 1680.
\item \textit{JKC3H8 v.\ Colton} (2013) 221 Cal.App.4th 468; \textit{Nguyen-Lam v.\ Cao} (2009) 171 Cal.App.4th 871.
\item \textit{Egan v.\ Mutual of Omaha Ins.\ Co.} (1979) 24 Cal.3d 809; \textit{Wilson v.\ 21st Century Ins.\ Co.} (2007) 42 Cal.4th 713.
\item \textit{Engalla v.\ Permanente Medical Group} (1997) 15 Cal.4th 951.
\end{enumerate}

\input{SIGBLOCK.tex}
\end{document}
